%! Author = lili
%! Date = 9/7/26

\section{Einleitung}\label{sec:einleitung}

Es gibt mehrere Dokumente, aus denen man lernen sollte:

\begin{itemize}
    \item Das Skript zum Praktikum
    \begin{itemize}
        \item Dazu gibt es Notizen aus dem Praktikum, die sich lohnen zu lesen
    \end{itemize}
    \item Versuchsfolien
    \begin{itemize}
        \item Je nach Versuch 1--3 PDFs
    \end{itemize}
\end{itemize}


\subsection{Was macht einen Modellorganismus aus?}\label{subsec:was-macht-einen-modellorganismus-aus?}

\begin{itemize}
    \item einfache Kultivierung unter Laborbedingun\gls{gen}
    \item kurze Generationszeit
    \item \gls{gen}etisch manipulierbar
    \item bei Eukaryoten: diploid, wenig repetitives (kompaktes) \gls{Genom}
    \item erlaubt Zugang zu wissenschaftlich oder ökonomisch wichtigen Fragestellun\gls{gen}
\end{itemize}

